\documentclass[12pt]{article}
\usepackage{amsmath, amssymb}
\usepackage[margin=1in]{geometry}
\setlength{\parskip}{6pt}
\title{QUBO Formulation: HP Lattice Protein Conformation Selection}
\date{}
\begin{document}
\maketitle

\subsection*{HP Lattice Protein Conformation Selection}

\paragraph{Problem Statement.}
Given a hydrophobic-polar protein sequence and a curated set of $M$ candidate lattice conformations, each associated with a precomputed contact energy $E_j$, the task is to select exactly one conformation that minimizes the total contact energy. The contact energy of a conformation is determined by the number of nonconsecutive hydrophobic-hydrophobic contacts it forms, with each contact contributing $-1$ to the total energy.

\paragraph{Decision Variables.}
For each candidate conformation $j \in \{0, 1, \dots, M-1\}$, we introduce a binary decision variable $x_j \in \{0, 1\}$. The variable $x_j$ equals $1$ if conformation $j$ is selected, and $0$ otherwise. The vector of decision variables is denoted $\mathbf{x} = (x_0, x_1, \dots, x_{M-1})^\top$.

\paragraph{QUBO Derivation.}
We construct the Quadratic Unconstrained Binary Optimization (QUBO) formulation by converting the constrained minimization problem into an unconstrained energy minimization problem via a penalty method.

\textbf{Step 1.} The original objective is to minimize the total contact energy over the binary variables:
\begin{equation}
\min_{\mathbf{x} \in \{0,1\}^M} \sum_{j=0}^{M-1} E_j x_j.
\end{equation}

\textbf{Step 2.} The selection requirement is enforced by the equality constraint that exactly one conformation must be chosen:
\begin{equation}
\sum_{j=0}^{M-1} x_j = 1.
\end{equation}

\textbf{Step 3.} We convert the constraint into a quadratic penalty term with weight $\lambda > 0$:
\begin{align}
P(\mathbf{x}) &= \lambda \left( \sum_{j=0}^{M-1} x_j - 1 \right)^2 \\
&= \lambda \left( \sum_{j=0}^{M-1} x_j^2 + \sum_{j \neq k} x_j x_k - 2 \sum_{j=0}^{M-1} x_j + 1 \right) \\
&= \lambda \left( \sum_{j=0}^{M-1} x_j + \sum_{j \neq k} x_j x_k - 2 \sum_{j=0}^{M-1} x_j + 1 \right) \\
&= \lambda \left( \sum_{j \neq k} x_j x_k - \sum_{j=0}^{M-1} x_j + 1 \right),
\end{align}
where the second equality follows from expanding the square, and the third equality uses the idempotent property $x_j^2 = x_j$ for binary variables. The cross-term sum $\sum_{j \neq k} x_j x_k$ can be rewritten over unique pairs as $2 \sum_{0 \leq j < k \leq M-1} x_j x_k$.

\textbf{Step 4.} Combining the objective and the penalty term yields the QUBO Hamiltonian $H(\mathbf{x})$:
\begin{align}
H(\mathbf{x}) &= \sum_{j=0}^{M-1} E_j x_j + \lambda \left( 2 \sum_{0 \leq j < k \leq M-1} x_j x_k - \sum_{j=0}^{M-1} x_j + 1 \right) \\
&= \sum_{j=0}^{M-1} (E_j - \lambda) x_j + 2\lambda \sum_{0 \leq j < k \leq M-1} x_j x_k + \lambda.
\end{align}

\textbf{Step 5.} Reading off the coefficients from the standard QUBO form $H(\mathbf{x}) = \sum_{j \leq k} Q_{jk} x_j x_k + c$, we obtain the general closed-form expressions for the upper-triangular Q matrix entries and the constant offset:
\begin{align}
Q_{jj} &= E_j - \lambda, \quad \forall j \in \{0, \dots, M-1\}, \\
Q_{jk} &= 2\lambda, \quad \forall 0 \leq j < k \leq M-1, \\
c &= \lambda.
\end{align}

\paragraph{Penalty Weight Justification.}
The penalty weight $\lambda$ must be chosen sufficiently large to ensure that any constraint violation is energetically unfavorable compared to the best feasible solution. The minimum penalty incurred by an infeasible assignment (where $\sum_{j=0}^{M-1} x_j \neq 1$) is $\lambda$, occurring when the sum equals $0$ or $2$. The maximum possible reduction in the objective energy achievable by switching from a feasible state to an infeasible state is bounded by $\max_{j} |E_j|$. Therefore, setting
\begin{equation}
\lambda > \max_{j \in \{0, \dots, M-1\}} |E_j|
\end{equation}
guarantees that the penalty cost strictly dominates any potential objective gain from violating the selection constraint.

\paragraph{Correctness Argument.}
Minimizing $H(\mathbf{x})$ over $\{0,1\}^M$ recovers the optimal conformation because feasible solutions satisfy $\sum_{j=0}^{M-1} x_j = 1$ and thus incur zero penalty, reducing $H(\mathbf{x})$ exactly to the original objective $\sum_{j=0}^{M-1} E_j x_j$. Infeasible solutions incur a penalty of at least $\lambda$. Since $\lambda$ exceeds the maximum possible objective value magnitude, any infeasible assignment has strictly higher energy than the optimal feasible assignment. Consequently, the global minimum of $H(\mathbf{x})$ must be feasible and correspond to the conformation with the lowest contact energy.

\paragraph{Illustrative Example.}
Consider a small instance with $M=3$ candidate conformations and precomputed energies $E_0$, $E_1$, and $E_2$. Substituting these parameters into the general formulas yields the concrete Hamiltonian:
\begin{align}
H(\mathbf{x}) &= E_0 x_0 + E_1 x_1 + E_2 x_2 + 2\lambda(x_0 x_1 + x_0 x_2 + x_1 x_2) - \lambda(x_0 + x_1 + x_2) + \lambda.
\end{align}
The corresponding Q matrix entries and offset are:
\begin{align}
Q_{00} &= E_0 - \lambda, \quad Q_{11} = E_1 - \lambda, \quad Q_{22} = E_2 - \lambda, \\
Q_{01} &= Q_{02} = Q_{12} = 2\lambda, \quad c = \lambda.
\end{align}
With $\lambda > \max(|E_0|, |E_1|, |E_2|)$, the global minimum of $H(\mathbf{x})$ occurs at the binary vector $\mathbf{x}^*$ where $x_j^* = 1$ for the index $j$ that minimizes $E_j$, and $x_k^* = 0$ for all $k \neq j$. The minimum energy value is exactly $\min(E_0, E_1, E_2)$, correctly identifying the optimal conformation and its contact energy.

\end{document}